\documentclass{article}

% --- Required Packages ---
\usepackage[a4paper, margin=0in]{geometry}
\usepackage{fontspec}
\usepackage[svgnames]{xcolor}
\usepackage{graphicx}
\usepackage{tikz}
\usetikzlibrary{positioning}
\usepackage{hyperref}
\hypersetup{
    colorlinks=true,
    linkcolor=blue,
    filecolor=magenta,      
    urlcolor=cyan,
    pdftitle={Krider's Field Guide}, % Corrected document title
    pdfpagemode=FullScreen,
    }

% --- COLOR PALETTE ---
\definecolor{PaperColor}{HTML}{FDF5E6}
\definecolor{TextColor}{HTML}{36454F}
\definecolor{KriderYellow}{HTML}{FDF5E6}
\definecolor{KriderGreen}{HTML}{5E7756}
\definecolor{KriderDarkText}{HTML}{36454F}

% --- Document-wide Settings ---
\setmainfont{Roboto}[
    UprightFont = *,
    BoldFont = * Bold,
    ItalicFont = * Italic
]
\pagenumbering{gobble}

% --- Document Body ---
\begin{document}
\pagestyle{empty}

%---------------------------------------------------
% PAGE 1: Krider's Field Guide Cover (Corrected Layout)
%---------------------------------------------------
\begin{tikzpicture}[remember picture, overlay]
    \coordinate (C) at (current page.center);

    % --- Background ---
    \fill[KriderYellow] (current page.south west) rectangle (current page.south east);

    % --- Header Block ---
    \node[KriderDarkText, anchor=north west, align=left] at (current page.north west) [xshift=1.5cm, yshift=-1.5cm] {
        \normalfont \fontsize{10}{12}\selectfont Powered by \{data\}syntax
    };
    
    % --- Main Title (Above Image) ---
    \node[KriderDarkText, anchor=center, text centered] at (current page.north) [yshift=-3cm] {
        \bfseries \fontsize{50}{66}\selectfont Krider's  Field Guide
    };
    
    % --- Subheadings ---
    \node[KriderDarkText, anchor=center] at (current page.north) [yshift=-5cm] {
        \normalfont\itshape \fontsize{16}{20}\selectfont Continuing a Legacy of Horticultural Excellence Since 1906 
    };

    % --- Main Image (Resized and placed below the title) ---
    \node[anchor=center] at ([yshift=0.5cm]C) {
        \includegraphics[width=10.5cm]{cover.png}
    };
    
    % --- Bottom Product Box (with green background) ---
    \node[fill=KriderGreen, text=white, anchor=south, minimum width=\paperwidth-2cm, inner sep=15pt, outer sep=0, rounded corners=3pt] at (current page.south) [yshift=3.8cm] {
        \begin{minipage}{0.9\textwidth}
            \bfseries \fontsize{18}{20}\selectfont Garden Cosmos (Cosmos bipinnatus), the 'Rubenza' cultivar. \\[3mm]
            \normalfont \fontsize{11}{13}\selectfont (Illustrated Above). Commonly known as the Garden Cosmos or Mexican Aster. The 'Rubenza' cultivar has 3-inch blooms  that open as a deep, ruby red and then mature to a softer, antique rose, creating a multi-toned effect on a single plant. The flowers feature a bright yellow center, stand on tall, slender stems, and bloom from summer until the first frost. They are easy to grow, thrive in full sun, and are excellent for cutting gardens and attracting pollinators.
        \end{minipage}
    };
    
    % --- Krider Logo Footer (Bottom Left) ---
    \node[anchor=south west, align=left, text width=4in] at (current page.south west) [xshift=1.5cm, yshift=1.5cm] {
        \bfseries\itshape \fontsize{24}{24}\selectfont Krider's Field Guide \\
        \itshape\bfseries \fontsize{10}{12}\selectfont Week 43 - Fall - 2025 \\
        \normalfont\itshape \fontsize{9}{11}\selectfont Guide: 002
    };
    
    % --- Page Footer ---
    \node[KriderDarkText, anchor=center] at (current page.south) [yshift=20] {
        \fontsize{12}{14}\itshape Page 1 of 3
    };
\end{tikzpicture}

\newpage

%---------------------------------------------------
% PAGE 2: Professional 2x2 Rose Grid Page
%---------------------------------------------------
\begin{tikzpicture}[remember picture, overlay]
    \fill[KriderYellow] (current page.south west) rectangle (current page.north east);
    \coordinate (C) at (current page.center);

    % --- Main Title ---
    \node[KriderDarkText, anchor=north, text width=0.9\paperwidth, align=center] at (current page.north) [yshift=-3cm] {
        \bfseries\fontsize{36}{40}\selectfont 
    };

    % --- Define a style for each content block for consistency ---
    \tikzset{rose_block/.style={text width=8cm, align=center}}

    % --- Top Left Rose: Daisy ---
    \node[rose_block] at ([xshift=-5cm, yshift=6.5cm]C) {
        \includegraphics[width=7cm]{daisy.png} \\
        \vspace{5mm}
        \bfseries\fontsize{16}{18}\selectfont Ox-eye Daisy \\
        \normalfont\itshape\fontsize{11}{13}\selectfont(Leucanthemum vulgare) \\
        \normalfont\fontsize{11}{13}\selectfont The Ox-eye Daisy is a widespread \\
        perennial herb native to Europe and Asia, recognized by its flowerhead featuring 15-40 white ray florets surrounding a bright yellow center. \\
        
    };

    % --- Top Right Rose: Buccaneer ---
    \node[rose_block] at ([xshift=5cm, yshift=6.5cm]C) {
        \includegraphics[width=7cm]{redRose.png} \\
        \vspace{5mm}
        \bfseries\fontsize{16}{18}\selectfont Wichura's Rose \\
        \normalfont\itshape\fontsize{11}{13}\selectfont (Rosa lucieae) \\
        \normalfont\fontsize{11}{13}\selectfont The wild Rosa lucieae is a white, five-petaled rambling rose, but it is a key parent to many hardy climbing hybrids, such as this double-flowered red variety. \\
    };

    % --- Bottom Left Rose: Roundelay ---
    \node[rose_block] at ([xshift=-5cm, yshift=-5.5cm]C) {
        \includegraphics[width=7cm]{orangeRose.png} \\
        \vspace{5mm}
        \bfseries\fontsize{16}{18}\selectfont Orange Hybrid Rose \\
        \normalfont\fontsize{11}{13}\selectfont This modern hybrid rose is distinguished by its vibrant, scarlet-orange petals and large, double-flowered blooms. \\
        It displays the classic, high-centered spiral form prized in cutting roses, with numerous petals unfurling from a tight center. \\
    };

    % --- Bottom Right Rose: Montezuma ---
    \node[rose_block] at ([xshift=5cm, yshift=-5.5cm]C) {
        \includegraphics[width=7cm]{flower.png} \\
        \vspace{5mm}
        \bfseries\fontsize{16}{18}\selectfont Meadow Knapweed \\
        \normalfont\itshape\fontsize{11}{13}\selectfont (Centaurea × gerstlaueri) \\
        \normalfont\fontsize{11}{13}\selectfont Recognized by its thistle-like purple flowerheads. It is identified by the distinctive light-to-dark-brown, papery, and fringed bracts at the base of the flower, which appear intermediate between its two parent species. \\
    };

    % --- Page Footer ---
    \node[KriderDarkText, anchor=center] at (current page.south) [yshift=50] {
        \fontsize{12}{14}\itshape Page 2 of 3
    };
\end{tikzpicture}

\newpage

%---------------------------------------------------
% PAGE 3: CITATION PAGE (Themed)
%---------------------------------------------------
\begin{tikzpicture}[remember picture, overlay]
    \fill[KriderYellow] (current page.south west) rectangle (current page.north east);
    
    % --- Main Title ---
    \node[KriderDarkText, anchor=center] at (current page.center) [yshift=250] {
        \fontsize{32}{36}\bfseries Project Source \& Collaboration
    };
    
    % --- Explanatory Text ---
    \node[KriderDarkText, anchor=north, text width=0.7\paperwidth, align=center] at (current page.center) [yshift=180] {
        \fontsize{14}{22}\selectfont
        This document was written in \LaTeX. The complete source code, image assets, and project history are available on GitHub. Suggestions, and flower requests are welcome!.
    };

    % --- GitHub Logo and Link ---
    \node[anchor=center] at (current page.center) [yshift=-20] {
        \href{https://github.com/Data-Syntax/KridersFieldGuide}{\includegraphics[height=5cm]{github-mark.png}}
    };
    
    \node[KriderDarkText, anchor=center] at (current page.center) [yshift=-100] {
        \fontsize{14}{18}\selectfont\url{https://github.com/Data-Syntax/KridersFieldGuide}
    };
    
    % --- Page Footer ---
    \node[KriderDarkText, anchor=center] at (current page.south) [yshift=50] {
        \fontsize{12}{14}\itshape Page 3 of 3
    };
\end{tikzpicture}

\end{document}